% =====================================================================
%  Curriculum Vitae — Chang Hwan Kim
%  Compiles with pdfLaTeX or XeLaTeX (Overleaf / TeX Live full).
%  Edit the content below; layout macros are defined in the preamble.
% =====================================================================
\documentclass[10pt,letterpaper]{article}

% ---------- Packages ----------
\usepackage{fontawesome5}
\usepackage[T1]{fontenc}
\usepackage{lmodern}
\usepackage[letterpaper,top=0.65in,bottom=0.8in,left=0.75in,right=0.75in,footskip=0.35in]{geometry}
\usepackage[dvipsnames]{xcolor}
\usepackage{titlesec}
\usepackage{enumitem}
\usepackage{fancyhdr}
\usepackage{lastpage}
\usepackage{tabularx}
\usepackage[hidelinks]{hyperref}

% ---------- Colors ----------
\definecolor{linkblue}{RGB}{0,79,144}
\definecolor{softgray}{gray}{0.55}

% ---------- Basic settings ----------
\pagestyle{fancy}
\fancyhf{}
\renewcommand{\headrulewidth}{0pt}
\cfoot{\color{softgray}\small\textit{Chang Hwan Kim -- Page \thepage\ of \pageref*{LastPage}}}
\setlength{\parindent}{0pt}
\setlength{\parskip}{0pt}
\linespread{1.0}
\raggedright
\hypersetup{pdftitle={Chang Hwan Kim's CV},pdfauthor={Chang Hwan Kim},colorlinks=false}

% ---------- Section headings ----------
\titleformat{\section}{\large\bfseries}{}{0pt}{}[\vspace{-7pt}\rule{\textwidth}{0.6pt}]
\titlespacing*{\section}{0pt}{8pt}{1pt}
\titleformat{\subsection}{\normalsize\bfseries}{}{0pt}{}
\titlespacing*{\subsection}{0pt}{3pt}{2pt}

% ---------- Lists ----------
\setlist[itemize]{label=$\circ$,leftmargin=1.6em,labelsep=0.5em,itemsep=0pt,topsep=1pt,parsep=0pt,partopsep=0pt}

% ---------- Macros ----------
% Linked label with external-link icon, e.g. \extlink{URL}{PDF}
\newcommand{\extlink}[2]{\href{#1}{\color{linkblue}[#2]~\scriptsize\faExternalLink*}}
% Header contact item
\newcommand{\contact}[3]{\href{#2}{#1~#3}}
% Two-line entry: {bold left}{italic right}{italic left}{right}
\newcommand{\entry}[4]{\textbf{#1}\hfill\textit{#2}\par\nopagebreak\textit{#3}\hfill{#4}\par\nopagebreak}
% One-line entry: {bold left}{italic right}
\newcommand{\entryline}[2]{\textbf{#1}\hfill\textit{#2}\par\nopagebreak}
% Sub-line inside an entry: {italic left}{right}
\newcommand{\subline}[2]{\textit{#1}\hfill{#2}\par\nopagebreak}
% Project heading inside an entry: {title}{dates}
\newcommand{\project}[2]{\vspace{2pt}\begin{tabularx}{\textwidth}{@{}>{\raggedright\arraybackslash}X r@{}}\textbf{\textit{#1}} & \textit{#2}\end{tabularx}\par\nopagebreak}
% Gap between entries
\newcommand{\entrygap}{\vspace{3pt}}
% Publication item: {title}{authors}{venue and links}
\newcommand{\pub}[4]{\item[#1.] #2\\ #3\\ #4}
% Self name in author lists
\newcommand{\me}{\textbf{C.H. Kim}}

% Publication lists with prefixed labels (J1., C1., ...)
\newlist{publist}{enumerate}{1}
\setlist[publist]{leftmargin=2.4em,labelsep=0.6em,align=right,itemsep=4pt,topsep=2pt,parsep=0pt}

\begin{document}

% =====================================================================
%  Header
% =====================================================================
\vspace*{-6pt}
\hfill{\color{softgray}\small\textit{Last updated in Nov. 2026}}\par
\vspace{2pt}
\begin{center}
{\fontsize{24pt}{28pt}\selectfont\bfseries Chang Hwan Kim}\\[7pt]
\small
\contact{\faEnvelope[regular]}{mailto:hwankim@umich.edu}{hwankim@umich.edu}\hspace{1.4em}
\contact{\faGlobe}{https://chang-hwan-kim1.github.io/}{Website}\hspace{1.4em}
\contact{\faGraduationCap}{https://scholar.google.com/citations?user=DUZ1KAoAAAAJ}{Scholar}\hspace{1.4em}
\contact{\faLinkedinIn}{https://www.linkedin.com/in/chang-hwan-kim-24aa1b316/}{LinkedIn}\hspace{1.4em}
\contact{\faGithub}{https://github.com/chang-hwan-kim1}{GitHub}
\end{center}
\vspace{-8pt}

% =====================================================================
\section{Education}
% =====================================================================
\entry{University of Michigan}{Aug. 2026 -- Present}{Ph.D. in Mechanical Engineering}{Ann Arbor, MI, USA}
\begin{itemize}
  \item Advisor: Prof. Gabor Orosz
\end{itemize}
\entrygap
\entry{Yonsei University}{Sep. 2024 -- Aug. 2026}{M.S. in Mechanical Engineering}{Seoul, Korea}
\begin{itemize}
  \item \textbf{GPA}: 4.23/4.30 (3.97/4.00)
  \item Advisor: Prof. Jun Young Yoon
\end{itemize}
\entrygap
\subline{B.S. in Mechanical Engineering}{\textit{Mar. 2019 -- Aug. 2024}}
\begin{itemize}
  \item \textbf{GPA}: 4.13/4.30 (3.92/4.00), ranked 1st out of 145 in major, \textbf{Summa Cum Laude}
  \item Includes a 1.5-year leave of absence for mandatory military service (Jan. 2021 -- Jul. 2022)
\end{itemize}

% =====================================================================
\section{Research Interests}
% =====================================================================
Vehicle dynamics and control, model predictive control (MPC), control barrier functions (CBF), design and motion control of mechatronic systems

% =====================================================================
\section{Publications}
% =====================================================================
\subsection{Journal Articles}
\begin{publist}
  \pub{J5}{DNN-Approximated Offset-Free MPC on FPGA for 10-kHz Constrained In-Plane Control of 6-DOF Magnetic Levitation System}
      {\me, J.J. Kim, H.M. Yoon, J.Y. Yoon}
      {\textit{IEEE Transactions on Industrial Informatics} (\textbf{TII}), Early Access, 2026}
  \pub{J4}{Design and Control of Dual Reluctance-Lorentz 6-DOF Magnetic Levitator for In-Plane Long-Stroke Motion and Low-Power Vertical Levitation}
      {\me, H.M. Yoon, J.W. Jung, E.K. Kim, J.Y. Yoon}
      {\textit{IEEE/ASME Transactions on Mechatronics} (\textbf{TMECH}), under review (minor revision)}
  \pub{J3}{Active 6-DOF Magnetic Levitator for Fully-Automated High-Throughput Large-Range Single-Mode Optical Fiber Alignment}
      {H.M. Yoon, \me, K.J. Kang, J.M. Sung, E.K. Kim, \ldots, J.Y. Yoon}
      {\textit{IEEE Transactions on Industrial Electronics} (\textbf{TIE}) \extlink{https://chang-hwan-kim1.github.io/assets/files/journal_paper_fiber.pdf}{PDF}}
  \pub{J2}{Identification and Reduction of Normal-Direction Force Ripples in Permanent-Magnet Linear Synchronous Motors}
      {S.M. Lee, Y.S. Kwon, J.H. Jung, K.J. Kang, \me, J.Y. Yoon}
      {\textit{IEEE Transactions on Industrial Electronics} (\textbf{TIE}) \extlink{https://chang-hwan-kim1.github.io/assets/files/journal_paper_pmlsm.pdf}{PDF}}
  \pub{J1}{Study on Characteristics of SPGD Algorithm in Coherent Beam Combining with Phase and Tip-Tilt Control Speed Differences}
      {H.M. Yoon, S.M. Lee, J.W. Jung, K. Lee, J.J. Jung, \me, J.Y. Yoon}
      {\textit{Journal of the Korean Society for Precision Engineering} (\textbf{JKSPE}) \extlink{https://chang-hwan-kim1.github.io/assets/files/journal_paper_cbc.pdf}{PDF}}
\end{publist}

\subsection{Conference Presentations}
\begin{publist}
  \pub{C4}{Control Barrier Function-Based Safety-Critical Control of 6-DOF Magnetic Levitation System}
      {\me, J. Kim, H.M. Yoon, J.Y. Yoon}
      {\textit{Korean Society for Precision Engineering Spring Conference} (\textbf{KSPE}), May 2026 \extlink{https://chang-hwan-kim1.github.io/assets/files/conference_abstract_cbf_2026.pdf}{Abstract}}
  \pub{C3}{Offset-Free Model Predictive Control of 6-DOF Magnetic Levitator for High-Precision Trajectory Tracking}
      {\me, J. Kim, H.M. Yoon, J.Y. Yoon}
      {\textit{Korean Society for Precision Engineering Autumn Conference} (\textbf{KSPE}), Nov. 2025 \extlink{https://chang-hwan-kim1.github.io/assets/files/conference_abstract_mpc_2025.pdf}{Abstract}}
  \pub{C2}{Gain-Scheduled Feedforward-Feedback Control for Reducing Lateral Acceleration and Jerk of Vehicle During Lane Changing Maneuver}
      {\me, J. Kee, H. Jang, J. Wang, J.Y. Yoon}
      {\textit{International Conference on Precision Engineering and Sustainable Manufacturing} (\textbf{PRESM 2025}), Jul. 2025 \extlink{https://presm.org/data/Regular\%20Session/PP3-012.pdf}{Abstract}}
  \pub{C1}{Control Method for Maintaining Dynamic Performance of a 6-DOF Hybrid-Reluctance-Type Magnetic Levitator During Payload Mass Changes}
      {\me, H.M. Yoon, J. Jung, J.Y. Yoon}
      {\textit{International Conference on Precision Engineering and Sustainable Manufacturing} (\textbf{PRESM 2024}), Jul. 2024, \textbf{Outstanding Presentation Award} \extlink{https://chang-hwan-kim1.github.io/assets/files/conference_abstract_payload_2024.pdf}{Abstract}}
\end{publist}

% =====================================================================
\section{Research Experience}
% =====================================================================
\entryline{Yonsei University (Advisor: Prof. Jun Young Yoon)}{Seoul, Korea}
\subline{Graduate Researcher}{\textit{Sep. 2024 -- Aug. 2026}}
\subline{Undergraduate Researcher}{\textit{Jan. 2023 -- Aug. 2024}}

\project{FPGA-Based DNN-Approximated Offset-Free MPC for 6-DOF Magnetic Levitator}{Mar. 2025 -- Aug. 2026}
\begin{itemize}
  \item Designed an offset-free model predictive controller (MPC) with a disturbance observer for constrained in-plane motion control (C3).
  \item Approximated the MPC law with a deep neural network (DNN) and implemented it on an FPGA, achieving constrained control at 10 kHz (J5).
  \item Developed control barrier function (CBF)-based safety-critical control for the same platform (C4).
\end{itemize}

\project{Design and Control of 6-DOF Magnetic Levitator for Long-Stroke Zero-Power Gravity Compensation}{Jan. 2024 -- Jun. 2025}
\begin{itemize}
  \item Designed and fabricated a 6-DOF magnetic levitator that combines reluctance forces for gravity compensation with Lorentz forces for long-stroke in-plane motion (J4).
  \item Developed and implemented 6-DOF feedback and feedforward controllers for precision motion control; the prototype levitates stably with only 0.498 W at rest.
\end{itemize}

\project{Vehicle Lateral Dynamics Control for Reducing Lateral Acceleration and Jerk During Lane Changes}{Sep. 2024 -- Feb. 2025}
\begin{itemize}
  \item Developed a gain-scheduled preview feedforward controller and integrated it with a lane-keeping LQR feedback controller (C2).
  \item Used a single-track vehicle model with Ackermann steering geometry and designed a Luenberger-PI observer for state estimation.
  \item Reduced jerk by up to 79.6\% and tracking error by up to 54.2\% relative to feedback-only control over 60--120 km/h in simulation.
\end{itemize}

\project{Force Ripple Reduction Control of Permanent-Magnet Linear Synchronous Motor (PMLSM)}{Mar. 2023 -- Mar. 2024}
\begin{itemize}
  \item Designed a DQ-axis current control method that reduces both normal- and tangential-direction force ripples, achieving 69\% and 87\% RMS reduction, respectively (J2).
  \item Proposed a method to identify the motor parameters that induce force ripples in iron-cored PMLSMs.
\end{itemize}

\project{Adaptive Control of Low-Power 6-DOF Hybrid-Reluctance-Type Magnetic Levitator}{Mar. 2023 -- Dec. 2023}
\begin{itemize}
  \item Designed a disturbance observer-based method to estimate the payload of the system.
  \item Implemented a controller that maintains the dynamic performance of the system under payload changes (C1, undergraduate thesis).
\end{itemize}

% =====================================================================
\section{Industry-Sponsored Projects}
% =====================================================================
\entry{Intelligent Controller Development for PAPS-Based Robotic C-Arm}{Sep. 2025 -- Aug. 2026}{In collaboration with GENORAY Co., Ltd.}{}
\begin{itemize}
  \item Designed sub-optimal collision-free paths and trajectories for a 4-DOF robotic C-arm using the A* algorithm.
  \item Implemented and tested the planners in a MATLAB simulation environment to ensure avoidance of the patient area.
\end{itemize}
\entrygap
\entry{Independent Drive and Control System for Styler Moving Hanger}{Nov. 2024 -- Oct. 2025}{In collaboration with LG Electronics Inc.}{}
\begin{itemize}
  \item Designed a novel drive mechanism enabling independent movement of the styler's moving hanger.
  \item Developed a motion controller that compensates for actuator cogging force and bearing friction.
\end{itemize}
\entrygap
\entry{Laser CBC SPGD Control Modeling and Simulation Development}{Apr. 2024 -- Mar. 2025}{In collaboration with LIG Nex1 Co., Ltd.}{}
\begin{itemize}
  \item Designed a stochastic parallel gradient descent (SPGD) algorithm for 19-channel laser coherent beam combining (J1).
  \item Developed a MATLAB simulator for laser beam modeling and SPGD performance evaluation.
\end{itemize}

% =====================================================================
\section{Honors and Scholarships}
% =====================================================================
\entryline{Summa Cum Laude (Top 1\%)}{Aug. 2024}
Graduated with highest honors from Yonsei University; one of four recipients in the College of Engineering.\par
\entrygap
\entryline{Outstanding Presentation Award}{Jul. 2024}
Awarded at the International Conference on Precision Engineering and Sustainable Manufacturing (PRESM 2024).\par
\entrygap
\entryline{Semester Honors, Yonsei University}{Spring 2019, Fall 2019, Fall 2023, Spring 2024}
\entrygap
\entryline{Accelerated B.S.--M.S. Program Scholarship, Yonsei University}{Sep. 2024 -- Aug. 2026}
Full tuition for 4 semesters, awarded for academic excellence upon early admission to the master's program.\par
\entrygap
\entryline{Brain Korea 21 (BK21) Scholarship, National Research Foundation of Korea}{Sep. 2024 -- Aug. 2026}
Awarded to support graduate research in strategic fields under Korea's national research initiative.\par
\entrygap
\entryline{Daesang Foundation Scholarship, Daesang Group}{Mar. 2020 -- Aug. 2024}
Full tuition for 6 semesters, awarded to outstanding undergraduate students.\par
\entrygap
\entryline{Yonsei Eminence Scholarship, Yonsei University}{Sep. 2019 -- Feb. 2020}
One-third tuition for one semester, awarded for academic excellence.\par

% =====================================================================
\section{Teaching Experience}
% =====================================================================
\entry{MEU3014 -- Mechatronics}{Fall 2024}{Teaching Assistant, Yonsei University}{}
\begin{itemize}
  \item Topics: electronics, sensors, motors, linear control theory, and control practice using a microcontroller.
\end{itemize}
\entrygap
\entry{MEU3007 -- Independent Study}{Spring 2025}{Teaching Assistant, Yonsei University}{}
\begin{itemize}
  \item Supervised an undergraduate student on their undergraduate thesis.
\end{itemize}

% =====================================================================
\section{Service and Other Experience}
% =====================================================================
\entry{Republic of Korea Army}{Jan. 2021 -- Jul. 2022}{Infantry Rifle Squad Leader}{Inje, Korea}
\begin{itemize}
  \item Served in the Opposing Force (OPFOR) regiment, 2nd Mountain Regiment, Korea Combat Training Center (KCTC).
\end{itemize}
\entrygap
\entry{Korean Red Cross Blood Management Center}{Jun. 2017 -- Present}{Regular Blood Donor (20 donations since Jun. 2017)}{Seoul, Korea}

% =====================================================================
\section{Skills}
% =====================================================================
\begin{tabularx}{\textwidth}{@{}l X@{}}
  \textbf{Programming} & MATLAB/Simulink, Python, C, LabVIEW, \LaTeX \\
  \textbf{Toolboxes} & MATLAB MPC, Control System, Optimization, Deep Learning, and Robotics Toolboxes; Vehicle Dynamics Blockset; MPT3 \\
  \textbf{Design \& Simulation} & ANSYS Electronics, Autodesk Inventor \\
  \textbf{Languages} & Korean (native), English (full professional proficiency) \\
\end{tabularx}

\end{document}
